Penyusunan laporan skripsi ini dibagi menjadi lima bab utama yang saling terhubung dengan sistematika sebagai berikut:

\noindent\textbf{BAB 1 PENDAHULUAN}
\\ Menguraikan latar belakang masalah mengenai beban akumulasi data historis mahasiswa di Pusdatin UNTAR, rumusan masalah, batasan masalah terkait kriteria dinamis dan arsitektur terdistribusi \textit{Client-Worker}, tujuan dan manfaat penelitian, metodologi SDLC yang digunakan, lokasi dan waktu pelaksanaan, serta sistematika penulisan.

\noindent\textbf{BAB 2 DASAR TEORETIS, METODOLOGI, PERANCANGAN}
\\ Menyajikan kajian pustaka dan dasar teori yang relevan meliputi konsep \textit{data archiving}, \textit{on-demand retrieval}, arsitektur ASP.NET Web Forms, komunikasi HTTP API JSON (\textit{agent.ashx}), pemrosesan skrip PowerShell, serta fitur automatisasi SQL Server Agent. Bab ini juga memuat perancangan arsitektur terdistribusi, perancangan basis data arsip, dan rancangan antarmuka portal \textit{web}.

\noindent\textbf{BAB 3 PEMBUATAN DAN IMPLEMENTASI}
\\ Membahas kebutuhan perangkat lunak dan keras pada lingkungan \textit{development/staging}, tahap pembuatan modul \textit{handler} API JSON, konfigurasi tugas penjadwalan pada SQL Server Agent dan PowerShell \textit{worker}, serta pembuatan halaman portal \textit{web} ASP.NET untuk pencarian data.

\noindent\textbf{BAB 4 PENGUJIAN DAN PEMBAHASAN}
\\ Menjelaskan hasil pengujian fungsionalitas dan performa sistem pada lingkungan pengembangan, analisis kecepatan pemindahan data historis, validasi pencarian data lintas basis data, serta pembahasan terhadap pencapaian solusi yang dirancang.

\noindent\textbf{BAB 5 KESIMPULAN DAN SARAN}
\\ Menyajikan kesimpulan akhir berdasarkan hasil pengujian dan pembahasan yang telah dilakukan, serta memberikan saran-saran akademis maupun praktis untuk pengembangan sistem pengarsipan di masa mendatang.